\documentclass[11pt]{article}

\usepackage[margin=1in]{geometry}
\usepackage[T1]{fontenc}
\usepackage{lmodern}
\usepackage{amsmath,amssymb,amsthm}
\usepackage{bm}
\usepackage{booktabs}
\usepackage{graphicx}
\usepackage{siunitx}
\usepackage[round,numbers,sort&compress]{natbib}
\usepackage[colorlinks=true,linkcolor=blue,citecolor=blue,urlcolor=blue]{hyperref}

\title{Background and related work}
\author{Sergio Lucero}
\date{\today}

\begin{document}
\maketitle

\section{Background and related work}
\label{chap:related}

This chapter places the thesis in four literatures it borrows from and tries to
satisfy at once: subsurface flow and transport simulation, learned surrogates for
partial differential equations, optimization under uncertainty for remediation
design, and the verification and statistical practice that decides whether a
computed number may be believed. The last of these is the one the thesis argues
about.

\subsection{Flow and transport in heterogeneous aquifers}

Steady groundwater flow in a confined aquifer is governed by Darcy's law with a
hydraulic conductivity that varies over several orders of magnitude in space.
Analytical solutions exist only in idealized configurations, and two of them are
used here as references: the Thiem solution for steady radial flow to a well
\citep{thiem1906,bear1972} and the Ogata--Banks solution for one-dimensional
longitudinal dispersion from a step source \citep{ogata1961solution}. Both are
old, both are derived under assumptions (homogeneity, one dimensionality, a fully
penetrating well) that a real site violates, and that is exactly why they are
useful: a numerical model that cannot reproduce them is not wrong about the site,
it is wrong about the equation.

Contaminant transport couples advection by the seepage velocity with hydrodynamic
dispersion and, where relevant, retardation and decay. The standard treatment of
heterogeneity is geostatistical: conductivity is modelled as a log-normal random
field with an exponential covariance, discretized here by a Karhunen--Lo\`eve
expansion truncated at $95\%$ of the field variance. The truncation is not a
modelling convenience but a dimension decision --- at the resolution used in this
thesis it retains $1\,529$ modes, which is the number any surrogate has to face.

Finite-volume discretizations with upwind advection \citep{leveque2002finite} are
the workhorse for this class of problem, and they are what the reference
implementation of \S\ref{chap:formulation} and the batched GPU backend both
implement, so that the difference between them is arithmetic order and data type,
not physics.

\subsection{Learned surrogates for partial differential equations}

Physics-informed neural networks \citep{raissi2019pinn,karniadakis2021piml}
embed the residual of the governing equation in the training loss, so that a
network can be trained without solution data. Operator-learning methods instead
learn a map between function spaces: DeepONet \citep{lu2021deeponet} from a
branch/trunk decomposition, and the Fourier neural operator
\citep{li2020fourier,li2021fno,kovachki2023neural} by multiplying retained Fourier
modes inside each layer, which gives it a global receptive field and a
resolution-independent parameterization.

The literature on failure modes is as important as the literature on successes.
Spectral bias makes networks learn low frequencies first and sharp fronts late;
gradient pathologies under stiff residuals and multi-objective losses are
documented and partly remedied by adaptive weighting
\citep{wang2021gradient,jagtap2020xpinn}; and the observation that PINN training
degrades as the problem becomes harder \citep{krishnapriyan2021failure} is a
warning against extrapolating from the easy cases.

A recurrent methodological hazard, and the one that this thesis spends a chapter
on, is that surrogate accuracy is usually reported as a single normalized error
for a single training run. Relative $L^2$ error has a built-in reference point ---
the zero predictor scores exactly $1.0$ --- which makes it possible to report a
number near $1.0$ without noticing that the surrogate is worse than useless. The
literature rarely reports the trivial-predictor baseline, the seed spread, or the
effect of the resolution the operator was trained at, and all three decide whether
a claimed improvement is real.

\subsection{Optimization under uncertainty for remediation design}

Simulation--optimization for pump-and-treat design has a long history
\citep{gorelick1983,wagner1995,mayer2002,ahlfeld2000}, reviewed in
\citep{zheng2002applied}. The classical formulations are deterministic: minimize
pumping cost subject to a concentration constraint evaluated on one calibrated
model. Because that model is uncertain, the design is fragile in exactly the way
the constraint is supposed to prevent, which motivates robust and
chance-constrained variants.

Two risk measures dominate. Value-at-risk constrains a quantile, is
non-convex and ignores the size of violations; conditional value-at-risk (CVaR)
averages the tail, is convex in the underlying distribution, and admits a linear
or smooth reformulation via the Rockafellar--Uryasev construction
\citep{rockafellar1991,rockafellar2000}. This thesis uses a CVaR penalty with
squared hinge, which keeps the objective smooth while making a policy that is
inside the constraint strictly better than one that merely touches it.

\subsection{Sample average approximation and what it certifies}

Chance-constrained and expectation objectives over a continuous random field are
not computable in closed form, so they are replaced by a sample average over $M$
realizations, an approach analysed in \citep{kleywegt2002,birge2011,shapiro2009,
nemirovski2009}. The theory says two things that matter for practice: the
optimal value of the sample problem converges to the true one as $M$ grows, and
the optimizer of the sample problem is \emph{not} certified by the sample it was
optimized on. A candidate solution has to be re-evaluated on independent draws,
and the resulting estimator of the reliability has its own sampling error.

That second step is often skipped, and the reason is arithmetic rather than
laziness: it costs as much again as the optimization. The consequence, measured
in \S\ref{chap:results}, is that a policy optimized on sixteen scenarios can
report a perfect reliability and fail on four hundred unseen ones, and that a
tail statistic computed on the optimization sample is not even approximately the
tail of the distribution.

\subsection{Derivative-free search and batched evaluation}

The objective defined by a transport solve is non-convex, non-smooth in the
constraints, and has no usable gradient in the presence of the random field; the
discrete adjoint exists but is not what the batched backend computes. The search
used here is therefore a compass (pattern) search with capacity projection, a
classical derivative-free method with convergence theory under mild smoothness
\citep{conn2009,nocedal2006}. Its cost structure is unusually friendly to
batching: a whole sweep of compass directions can be evaluated as one batch of
initial-value problems, which turns a sequence of small linear solves into one
large one and is where the GPU speedups of \S\ref{chap:results} come from.

\subsection{Verification, validation, and reproducible computational results}

The vocabulary used throughout is the standard one. \emph{Code verification}
asks whether the implementation solves the intended equations, and is answered
by convergence studies and manufactured solutions; \emph{solution verification}
asks whether a particular discretization is accurate enough; \emph{validation}
asks whether the equations describe the site. This thesis is concerned almost
entirely with the first, plus a fourth activity that the literature treats
informally: \emph{implementation parity}, i.e.\ whether two independently written
backends of the same model agree. Parity is the cheapest available check on a
fast implementation, and it is the one that catches the errors that benchmarks
miss.

The reproducibility literature \citep{sandve2013ten,wilson2014best} has
established the practices --- version control, tests, recorded environments,
archived artifacts --- and equally that they are not followed under deadline
pressure. The argument developed in this thesis is narrower and more operational:
each verification device should be tied to a defect it would have caught, and it
should be cheap enough to run on every change. Where a device is expensive (an
out-of-sample certificate, a seed sweep), the cost is made explicit so that the
decision to skip it is a decision and not an oversight.

\subsection{Statistical reporting of rare events}

A reliability of $0.975$ estimated from forty scenarios and the same number
estimated from four hundred are not the same claim, and a tail average is weaker
still. The interval used here is the Wilson score interval for a binomial
proportion \citep{wilson1927probable}, chosen over the normal approximation
because it behaves sensibly at the extremes that matter for a $5\%$ failure
probability. For the tail itself the relevant fact is that a CVaR estimate at
level $\alpha$ averages $\alpha M$ of the most variable draws: at
$\alpha = 0.05$ and $M = 400$ that is twenty draws, so its standard error is
$\operatorname{sd}(\text{tail})/\sqrt{20}$ --- of order a third of a unit for the
compliance metric used here. No amount of careful implementation removes that
uncertainty; the only responses are to report it or to change the sample size,
and \S\ref{chap:results} quantifies the trade-off.

\subsection{Positioning of this work}

The contributions of the thesis are not new numerical methods. The batched
transport solve, the CVaR formulation, the pattern search and the Wilson interval
are all standard. What is specific here is the \emph{protocol} that decides when a
number may be reported: seed sweeps instead of best runs for surrogates,
implementation parity as a gate on the fast backend, an out-of-sample
certificate with an interval for reliability, and an explicit standard error for
the tail. Applied consistently, that protocol produces two results that a
conventional workflow would have hidden: a surrogate that was worse than
predicting zeros, and a tail statistic that describes the sample it was computed
on rather than the aquifer it is supposed to describe.

\begin{thebibliography}{99}

\bibitem[Thiem(1906)]{thiem1906}
G.~Thiem.
\textit{Hydrologische Methoden}.
Gebhardt, Leipzig, 1906.

\bibitem[Bear(1972)]{bear1972}
J.~Bear.
\textit{Dynamics of Fluids in Porous Media}.
American Elsevier, New York, 1972.

\bibitem{ogata1961solution}
A. Ogata and R. B. Banks,
\textit{A solution of the differential equation of longitudinal dispersion in
porous media}, U.S. Geological Survey Professional Paper 411-A, 1961.

\bibitem[LeVeque(2002)]{leveque2002finite}
R.~J. LeVeque.
\textit{Finite Volume Methods for Hyperbolic Problems}.
Cambridge University Press, 2002.

\bibitem[Raissi et~al.(2019)]{raissi2019pinn}
M.~Raissi, P.~Perdikaris, and G.~E. Karniadakis.
Physics-informed neural networks: A deep learning framework for solving forward
and inverse problems involving nonlinear partial differential equations.
\textit{Journal of Computational Physics}, 378:686--707, 2019.

\bibitem[Karniadakis et~al.(2021)]{karniadakis2021piml}
G.~E. Karniadakis, I.~G. Kevrekidis, L.~Lu, P.~Perdikaris, S.~Wang, and L.~Yang.
Physics-informed machine learning.
\textit{Nature Reviews Physics}, 3:422--440, 2021.

\bibitem[Lu et~al.(2021)]{lu2021deeponet}
L.~Lu, P.~Jin, G.~Pang, Z.~Zhang, and G.~E. Karniadakis.
Learning nonlinear operators via DeepONet based on the universal approximation
theorem of operators.
\textit{Nature Machine Intelligence}, 3:218--229, 2021.

\bibitem{li2020fourier}
Z. Li, N. Kovachki, K. Azizzadenesheli, B. Liu, K. Bhattacharya, A. Stuart,
A. Anandkumar, \textit{Fourier Neural Operator for Parametric Partial
Differential Equations}, ICLR, 2021.

\bibitem[Li et~al.(2021)]{li2021fno}
Z.~Li, N.~Kovachki, K.~Azizzadenesheli, B.~Liu, K.~Bhattacharya, A.~Stuart,
and A.~Anandkumar.
Fourier neural operator for parametric partial differential equations.
\textit{International Conference on Learning Representations (ICLR)}, 2021.

\bibitem[Kovachki et~al.(2023)]{kovachki2023neural}
N.~Kovachki, Z.~Li, B.~Liu, K.~Azizzadenesheli, K.~Bhattacharya, A.~Stuart,
and A.~Anandkumar.
Neural operator: Learning maps between function spaces with applications to
PDEs.
\textit{Journal of Machine Learning Research}, 24(89):1--97, 2023.

\bibitem[Wang et~al.(2021)]{wang2021gradient}
S.~Wang, Y.~Teng, and P.~Perdikaris.
Understanding and mitigating gradient flow pathologies in physics-informed
neural networks.
\textit{SIAM Journal on Scientific Computing}, 43(5):A3055--A3081, 2021.

\bibitem[Jagtap et~al.(2020)]{jagtap2020xpinn}
A.~D. Jagtap, E.~Kharazmi, and G.~E. Karniadakis.
Conservative physics-informed neural networks on discrete domains for
conservation laws: Applications to forward and inverse problems.
\textit{Computer Methods in Applied Mechanics and Engineering},
365:113028, 2020.

\bibitem[Krishnapriyan et~al.(2021)]{krishnapriyan2021failure}
A.~S. Krishnapriyan, A.~Gholami, S.~Zhe, R.~Kirby, and M.~W. Mahoney.
Characterizing possible failure modes in physics-informed neural networks.
\textit{Advances in Neural Information Processing Systems}, 34, 2021.

\bibitem[Gorelick(1983)]{gorelick1983}
S.~M. Gorelick.
A review of distributed parameter groundwater management modeling methods.
\textit{Water Resources Research}, 19(2):305--319, 1983.

\bibitem[Wagner(1995)]{wagner1995}
B.~J. Wagner.
Recent advances in simulation-optimization groundwater management modeling.
\textit{Reviews of Geophysics}, 33(S2):1021--1028, 1995.

\bibitem[Mayer et~al.(2002)]{mayer2002}
A.~S. Mayer, C.~T. Kelley, and C.~T. Miller.
Optimal design of pump-and-treat systems using simulation-optimization.
\textit{Water Resources Research}, 38(12):1277, 2002.

\bibitem[Ahlfeld and Mulligan(2000)]{ahlfeld2000}
D.~P. Ahlfeld and A.~E. Mulligan.
\textit{Optimal Management of Groundwater Resources}.
Wiley, New York, 2000.

\bibitem[Zheng and Bennett(2002)]{zheng2002applied}
C.~Zheng and G.~D. Bennett.
\textit{Applied Contaminant Transport Modeling}.
Wiley, New York, 2nd edition, 2002.

\bibitem[Rockafellar and Wets(1991)]{rockafellar1991}
R.~T. Rockafellar and R.~J.-B. Wets.
Scenarios and policy aggregation in optimization under uncertainty.
\textit{Mathematics of Operations Research}, 16(1):119--147, 1991.

\bibitem[Rockafellar and Uryasev(2000)]{rockafellar2000}
R.~T. Rockafellar and S.~Uryasev.
Optimization of conditional value-at-risk.
\textit{Journal of Risk}, 2(3):21--41, 2000.

\bibitem[Kleywegt et~al.(2002)]{kleywegt2002}
A.~J. Kleywegt, A.~Shapiro, and T.~Homem-de-Mello.
The sample average approximation method for stochastic discrete optimization.
\textit{SIAM Journal on Optimization}, 12(2):479--502, 2002.

\bibitem[Birge and Louveaux(2011)]{birge2011}
J.~R. Birge and F.~Louveaux.
\textit{Introduction to Stochastic Programming}.
Springer, New York, 2nd edition, 2011.

\bibitem[Shapiro et~al.(2009)]{shapiro2009}
A.~Shapiro, D.~Dentcheva, and A.~Ruszczy\'nski.
\textit{Lectures on Stochastic Programming: Modeling and Theory}.
SIAM, Philadelphia, 2009.

\bibitem[Nemirovski et~al.(2009)]{nemirovski2009}
A.~Nemirovski, A.~Juditsky, G.~Lan, and A.~Shapiro.
Robust stochastic approximation approach to stochastic programming.
\textit{SIAM Journal on Optimization}, 19(4):1574--1609, 2009.

\bibitem[Conn et~al.(2009)]{conn2009}
A.~R. Conn, K.~Scheinberg, and L.~N. Vicente.
\textit{Introduction to Derivative-Free Optimization}.
SIAM, Philadelphia, 2009.

\bibitem[Nocedal and Wright(2006)]{nocedal2006}
J.~Nocedal and S.~J. Wright.
\textit{Numerical Optimization}.
Springer, New York, 2nd edition, 2006.

\bibitem[Sandve et al.(2013)]{sandve2013ten}
G.~K. Sandve, A.~Nekrutenko, J.~Taylor, and E.~Hovig.
\newblock Ten simple rules for reproducible computational research.
\newblock \emph{PLoS Computational Biology}, 9(10):e1003285, 2013.

\bibitem[Wilson et al.(2014)]{wilson2014best}
G.~Wilson, D.~A. Aruliah, C.~T. Brown, N.~P. Chue~Hong, M.~Davis, R.~T. Guy,
  S.~H. Haddock, K.~D. Huff, I.~M. Mitchell, M.~D. Plumbley, B.~Waugh,
  E.~P. White, and P.~Wilson.
\newblock Best practices for scientific computing.
\newblock \emph{PLoS Biology}, 12(1):e1001745, 2014.

\bibitem[Wilson(1927)]{wilson1927probable}
E.~B. Wilson.
\newblock Probable inference, the law of succession, and statistical inference.
\newblock \emph{Journal of the American Statistical Association},
  22(158):209--212, 1927.

\end{thebibliography}

\end{document}
